\section{Methodology}

We train a LoRA~\citep{lora} adapter on the defender $\mathcal{M}_d$ (rank 32, $\alpha_{\text{LoRA}}=64$, fp32, applied to all attention and MLP projections) using a multi-objective loss that balances \emph{representation divergence}---breaking similarity with a frozen anchor $\mathcal{M}_a$---against \emph{behavior preservation} on benign inputs. Hidden states are extracted at layer $\ell = \lfloor 0.5 \cdot L \rfloor$. The representation divergence component is built on Centered Kernel Alignment (CKA)~\citep{kornblith2019similarity}, which compares relational structure between models via Gram matrices and is therefore agnostic to hidden dimension (e.g., $d_d = 4096$ vs.\ $d_a = 2560$); see \cref{sec:related} for background. Given hidden states $\mathbf{X}_d, \mathbf{X}_a \in \R^{N \times d}$ at layer $\ell$:
\begin{equation}
    \text{CKA}(\mathbf{X}_d, \mathbf{X}_a) = \frac{\text{HSIC}(\tilde{\mathbf{K}}_d, \tilde{\mathbf{K}}_a)}{\sqrt{\text{HSIC}(\tilde{\mathbf{K}}_d, \tilde{\mathbf{K}}_d) \cdot \text{HSIC}(\tilde{\mathbf{K}}_a, \tilde{\mathbf{K}}_a)}}
\end{equation}
where $\tilde{\mathbf{K}} = \mathbf{H}\mathbf{K}\mathbf{H}$ are centered Gram matrices and $\text{HSIC}$~\citep{gretton2005hsic} is computed as $\text{HSIC}(\mathbf{A},\mathbf{B}) = \text{tr}(\mathbf{A}^\top\mathbf{B})$.
High CKA ($\approx 1$) means two models encode the same neighborhood structure---if an adversarial suffix navigates the anchor toward compliance, it follows a similar trajectory in the defender. Minimizing CKA reorganizes the defender's geometry so that the anchor's directions no longer align, and because CKA operates on $N \times N$ Gram matrices rather than $d$-dimensional vectors, this works across architectures with different hidden dimensions.

Our core defense loss applies this repulsion exclusively to harmful prompts, focusing disruption on the representation subspace that attacks exploit while leaving benign geometry intact (\cref{tab:cka_ablation}): $\loss_{\text{CKA}} = \text{CKA}(\mathbf{X}_d^{\text{adapted}}, \mathbf{X}_a)$. However, CKA repulsion alone---while sufficient for strong security---damages benign output quality at the strength required for full protection. We therefore combine it with four auxiliary objectives that allow reducing the repulsion weight $\gamma$ while maintaining both security and utility:
\begin{equation}
    \loss = \gamma \cdot \loss_{\text{CKA}} + \alpha \cdot \loss_{\text{refusal}} + \beta \cdot \loss_{\text{coherency}} + \varepsilon \cdot \loss_{\text{KL}} + \delta \cdot \loss_{\text{LM}}
\end{equation}
The refusal-direction loss ($\loss_{\text{refusal}}$, weight $\alpha$), inspired by the finding that refusal is mediated by a single direction in activation space~\citep{repeng, arditi2024refusal, li2023iti}, pushes harmful representations toward a pre-computed refusal direction $\mathbf{r}$ (the normalized difference between mean refusal and mean compliance hidden states): $\loss_{\text{refusal}} = -\frac{1}{N_h} \sum_{i \in \text{harmful}} \cos(\mathbf{h}_i^{\text{adapted}}, \mathbf{r})$. This provides targeted pressure on harmful representations, allowing $\gamma$ to be reduced to a level that preserves benign behavior. The coherency loss ($\loss_{\text{coherency}}$, weight $\beta$) anchors the adapted model to its base weights via MSE on hidden states: $\loss_{\text{coherency}} = \frac{1}{N} \sum_{i=1}^{N} w_i \cdot \norm{\mathbf{h}_i^{\text{adapted}} - \mathbf{h}_i^{\text{base}}}^2$, where $w_i$ is higher for benign prompts, strongly preserving benign representations while allowing the defense to modify harmful-region representations. The KL-divergence loss ($\loss_{\text{KL}}$, weight $\varepsilon$) complements this by preserving the output token distribution on benign and borderline prompts (including XSTest~\citep{xstest}). Finally, an optional LM loss ($\loss_{\text{LM}}$, weight $\delta$) adds cross-entropy supervision on refusal response tokens; this is needed only for Mistral ($\delta = 0.06$) due to its weaker base refusal behavior, and is set to zero for Qwen and Llama-3.

\paragraph{Existing defenses under cross-model transfer.}
To contextualize our approach, we test existing Mistral-based defenses---Circuit Breakers~\citep{zou2024circuitbreakers} and RepBend~\citep{repbend}---against GCG suffixes transferred from 20 source models (\cref{tab:baselines}). Both show residual ASR (1.2\% and 1.3\%, respectively), with Circuit Breakers producing garbled output rather than coherent refusals. Our defense reduces cross-model ASR to 0--5\% with competitive utility (\cref{tab:benchmark_results}).

\begin{table}[ht]
\centering
\caption{Mistral-7B defense ASR (\%) against cross-model GCG transfer (20 source models). CB = Circuit Breakers~\citep{zou2024circuitbreakers}, RB = RepBend~\citep{repbend}.}
\label{tab:baselines}
\setlength{\tabcolsep}{10pt}
\renewcommand{\arraystretch}{1.3}
\begin{tabular}{lccc}
\toprule
 & CB (M7) & RB (M7) & \textbf{Ours} \\
\midrule
\textbf{Self (mistral)}  & 0.0 & 0.0 & \textbf{1.0} \\
\textbf{Overall (20 sources)}  & 1.2 & 1.3 & \textbf{2.0} \\
\bottomrule
\end{tabular}
\setlength{\tabcolsep}{6pt}
\renewcommand{\arraystretch}{1.0}
\vspace{4pt}
\parbox{\linewidth}{\small
CB (M7): GraySwanAI/Mistral-7B-Instruct-v0.2-RR (Circuit Breakers).\\
RB (M7): AIM-Intelligence/RepBend\_Mistral-7B-Instruct-v0.2 (RepBend).\\
Ours: Mistral-7B-Instruct-v0.2 with our defense (lr\_l3\_hx\_g10: $\alpha\!=\!0.15$, $\gamma\!=\!1.0$, anchor: Llama-3-8B).}
\end{table}
